\begin{abstract}
Large language models trained on massive corpora of data from the web can memorize and reproduce sensitive or private data
raising both legal and ethical concerns.
Unlearning, or tuning models to forget information present in their training data, provides us with a way to protect private data after training.  
Although several methods exist for such unlearning, it is unclear to what extent they result in models equivalent to those where the data to be forgotten was never learned in the first place.
To address this challenge, we present \tofu{}, a Task of Fictitious Unlearning, as a benchmark aimed at helping deepen our understanding of unlearning.  
We offer a dataset of $200$ diverse synthetic author profiles, each consisting of 20 question-answer pairs, and a subset of these profiles called the \emph{forget set} that serves as the target for unlearning.
We compile a suite of metrics that work together to provide a holistic picture of unlearning efficacy.
Finally, we provide a set of baseline results from existing unlearning algorithms.
Importantly, none of the baselines we consider show effective unlearning motivating continued efforts to develop approaches for unlearning that effectively tune models so that they truly behave as if they were never trained on the forget data at all.
\end{abstract}